\documentclass[../../main.tex]{subfiles}

\begin{document}

\section{单调函数的连续性}\label{section:ch06-单调连续}

回忆一个函数定义为单调的, 若它是递增的或递减的. 单调函数在解决上面提出的问题上起了决定性的作用. 这有两个理由. 首先, Lebesgue 的一个定理 (见 6.2 节) 断言开区间上的单调函数是几乎处处可微的. 其次, Jordan 的一个定理 (见 6.3 节) 告诉我们闭有界区间上的非常一般的函数族, 即那些有界变差函数 (包括 Lipschitz 函数), 可以表示为两个单调函数的差, 因此它们在其定义域内部也是几乎处处可微的. 在这简短的预备的一节中我们考虑单调函数的连续性质.

\begin{theorem}\label{theorem:ch06-单调连续}
令 $f$ 为开区间 $(a,b)$ 上的单调函数. 则 $f$ 除了在 $(a,b)$ 中的可能的可数个点外连续.
\end{theorem}

\begin{proof}
假设 $f$ 是递增的. 此外, 假设 $(a,b)$ 是有界的且 $f$ 在闭区间 $[a,b]$ 上递增. 否则, 将 $(a,b)$ 表示为上升的开有界区间序列的并, 这些区间的闭包含于 $(a,b)$, 取该可数区间族的每个区间中不连续点的并集. 对每个 $x_{0}\in(a,b)$, $f$ 在点 $x_{0}$ 有左极限与右极限. 定义
\begin{align*}
f(x_{0}^{-})=\lim_{x\to x_{0}^{-}}f(x)=\sup\{f(x)\mid a<x<x_{0}\}
\end{align*}
与
\begin{align*}
f(x_{0}^{+})=\lim_{x\to x_{0}^{+}}f(x)=\inf\{f(x)\mid x_{0}<x<b\}.
\end{align*}
由于函数 $f$ 是递增的, $f(x_{0}^{-})\leqslant f(x_{0}^{+})$. 函数 $f$ 在 $x_{0}$ 不连续当且仅当 $f(x_{0}^{-})<f(x_{0}^{+})$, 在这种情形我们定义开 ``跳跃'' 区间 $J(x_{0})$ 为
\begin{align*}
J(x_{0})=\{y\mid f(x_{0}^{-})<y<f(x_{0}^{+})\}.
\end{align*}
每个跳跃区间包含于有界区间 $[f(a),f(b)]$, 而跳跃区间族是不交的. 因此, 对每个自然数 $n$, 仅存在有限个长度大于 $1/n$ 的跳跃区间. 于是 $f$ 的不连续点集是有限集的可数族的并, 因此是可数的.
\end{proof}

\begin{proposition}\label{proposition:ch06-不连续点集}
令 $C$ 为开区间 $(a,b)$ 的可数子集. 则存在 $(a,b)$ 上的递增函数, 该函数仅在 $(a,b)\sim C$ 中的点处连续.
\end{proposition}

\begin{proof}
若 $C$ 是有限的, 证明是显然的. 假设 $C$ 是可数无穷的. 令 $\{q_{n}\}_{n=1}^{\infty}$ 为 $C$ 的列举. 定义 $(a,b)$ 上的函数 $f$ 为
\begin{align*}
\text{对所有 }a<x<b,\quad f(x)=\sum_{\{n\mid q_{n}\leqslant x\}}\frac{1}{2^{n}}.
\end{align*}
由于公比小于 $1$ 的几何级数收敛, 所以 $f$ 是恰当定义的. 此外,
\begin{align}
\text{若 }a<u<v<b,\text{则 }f(v)-f(u)=\sum_{\{n\mid u<q_{n}\leqslant v\}}\frac{1}{2^{n}}.\label{align:ch06-跳跃和}
\end{align}
因此 $f$ 是递增的. 令 $x_{0}=q_{k}$ 属于 $C$. 则根据\eqref{align:ch06-跳跃和},
\begin{align*}
\text{对所有 }x<x_{0},\quad f(x_{0})-f(x)\geqslant\frac{1}{2^{k}}.
\end{align*}
因此 $f$ 在点 $x_{0}$ 不连续. 现在令 $x_{0}$ 属于 $(a,b)\sim C$. 令 $n$ 为自然数. 存在包含 $x_{0}$ 的开区间 $I$, 使得对 $1\leqslant k\leqslant n$, $q_{k}$ 不属于 $I$. 我们从\eqref{align:ch06-跳跃和}推出对所有 $x\in I$, $|f(x)-f(x_{0})|<1/2^{n}$. 因此 $f$ 在点 $x_{0}$ 连续.
\end{proof}

\begin{exercise}\label{exercise:ch06-01}
令 $C$ 为非退化闭有界区间 $[a,b]$ 的可数子集. 证明存在 $[a,b]$ 上的递增函数, 它仅在 $[a,b]\sim C$ 中的点处连续.
\end{exercise}

\begin{exercise}\label{exercise:ch06-02}
证明存在 $[0,1]$ 上的严格递增函数, 它仅在 $[0,1]$ 中的无理数处连续.
\end{exercise}

\begin{exercise}\label{exercise:ch06-03}
令 $f$ 为 $\mathbf{R}$ 的子集 $E$ 上的单调函数. 证明 $f$ 除了 $E$ 中的可能的可数个点外连续.
\end{exercise}

\begin{exercise}\label{exercise:ch06-04}
令 $E$ 为 $\mathbf{R}$ 的子集而 $C$ 是 $E$ 的可数子集. 是否存在 $E$ 上的单调函数, 它仅在 $E\sim C$ 的点处连续?
\end{exercise}

\end{document}